% Section: Related Work
% Outline (for review):
% - Attribution graphs and Circuit Tracer; Neuronpedia and autointerp
% - SAE/CLT and replacement models background
% - Geometric clustering baselines vs behavior-grounded grouping
% - Causal tracing context and limitations
% - Clarify novelty: automated, behavior-validated grouping and naming

\section{Related Work}
\label{sec:related}

\paragraph{Attribution graphs and platforms}
Attribution graphs formalize how token embeddings and features contribute to target logits~\cite{ameisen2025circuittracing}. Neuronpedia provides feature cards, autointerp labels, and graph tooling that enable interactive analysis at scale~\cite{neuronpedia2025}. For a comprehensive survey of circuits research methods and cross-organization replication efforts, see~\cite{lindsey2025landscape}.

\paragraph{SAE/CLT and replacement models}
Sparse Autoencoders and Cross-Layer Transcoders (CLT) expose sparse, often monosemantic features and linearized pathways for attribution~\cite{bricken2023monosemanticity,templeton2024scaling,crosscoders2024}. Replacement models enable computation of direct effects between features while freezing attention patterns and layer norms, yielding tractable causal attributions.

\paragraph{Geometric baselines vs behavioral grouping}
Clustering features using purely geometric information (e.g., cosine on activation vectors or adjacency in layer-influence space) can produce compact clusters but need not align with functional roles. Our approach deliberately prioritizes behavioral coherence derived from cross-prompt signatures, with quantitative comparisons to geometric baselines.

\paragraph{Causal tracing and interpretability}
Causal tracing techniques aim to identify mechanisms by intervention and controlled replacement. Our validation uses correlational graph metrics and descriptive behavior metrics; we outline planned ablations and steering experiments to strengthen causal claims in future work.

\paragraph{Novelty}
We automate concept hypothesis generation, cross-prompt measurement, and transparent rule-based grouping and naming. The novelty lies in leveraging behavior-grounded signals for grouping and labeling, not in proposing a new distance metric or clustering algorithm. Conceptually, our probe-prompt generator descends from Prompt Rover's ``semantic navigation'' framing, which we previously introduced as a black-box exploratory tool; here we formalize it into a rule-driven, pre-specified prompt family evaluated on attribution-graph features~\cite{birardi2025insight,birardi2025promptrover}.


